% ─── MÓDULO 11: Almacenamiento de datos relacionados ─────
\newpage
\section{Almacenamiento de Datos Relacionados}

Cuando una base de datos posee \textbf{integridad referencial} activa, el orden
y la forma en que se insertan los registros importan tanto como el contenido
mismo. Violar el orden de inserción provoca errores de clave foránea que el
motor rechaza antes de escribir cualquier dato.

\begin{nota}
  Regla general: insertar siempre primero la tabla \textbf{padre} (la que
  contiene la clave primaria referenciada) y después la tabla \textbf{hijo}
  (la que contiene la clave foránea). Para eliminar, el orden es el inverso.
\end{nota}

% ─────────────────────────────────────────────────────────
\subsection{Inserción en Tablas con Relación 1:1}

\begin{definicion}{Relación 1:1}
  Una fila de la tabla \textbf{A} se corresponde con \textbf{exactamente una}
  fila de la tabla \textbf{B}, y viceversa. Se implementa colocando la clave
  foránea en la tabla secundaria con restricción \texttt{UNIQUE}.
\end{definicion}

\textbf{Esquema de ejemplo:}

\begin{lstlisting}[style=sql, caption={Tablas con relación 1:1 — empleado y expediente}]
CREATE TABLE empleados (
    empleado_id   INT          UNSIGNED AUTO_INCREMENT PRIMARY KEY,
    nombre        VARCHAR(100) NOT NULL,
    departamento  VARCHAR(60)  NOT NULL
);

CREATE TABLE expedientes (
    expediente_id INT          UNSIGNED AUTO_INCREMENT PRIMARY KEY,
    empleado_id   INT          UNSIGNED NOT NULL UNIQUE,   -- UNIQUE fuerza 1:1
    nss           CHAR(11)     NOT NULL,
    fecha_alta    DATE         NOT NULL,
    FOREIGN KEY (empleado_id) REFERENCES empleados(empleado_id)
        ON DELETE CASCADE
        ON UPDATE CASCADE
);
\end{lstlisting}

\textbf{Orden correcto de inserción:}

\begin{lstlisting}[style=sql, caption={Inserción respetando integridad referencial 1:1}]
-- Paso 1: insertar en la tabla PADRE
INSERT INTO empleados (nombre, departamento)
VALUES ('Laura Mendoza', 'Ingeniería'),
       ('Roberto Solis', 'Finanzas');

-- Paso 2: recuperar los IDs generados (o usar LAST_INSERT_ID() para uno solo)
-- empleado_id = 1 -> Laura Mendoza
-- empleado_id = 2 -> Roberto Solis

-- Paso 3: insertar en la tabla HIJO referenciando los IDs existentes
INSERT INTO expedientes (empleado_id, nss, fecha_alta)
VALUES (1, '578-12-3456', '2021-03-15'),
       (2, '412-87-9900', '2019-07-01');
\end{lstlisting}

\begin{ejemplo}
  Uso de \texttt{LAST\_INSERT\_ID()} para encadenar inserciones individuales
  de forma segura:

\begin{lstlisting}[style=sql, caption={Encadenamiento con LAST\_INSERT\_ID()}]
INSERT INTO empleados (nombre, departamento)
VALUES ('Carmen Vidal', 'Diseño');

-- LAST_INSERT_ID() devuelve el AUTO_INCREMENT de la última inserción
INSERT INTO expedientes (empleado_id, nss, fecha_alta)
VALUES (LAST_INSERT_ID(), '301-55-7812', '2023-01-10');
\end{lstlisting}
\end{ejemplo}

\begin{nota}
  \texttt{LAST\_INSERT\_ID()} es \textbf{seguro en entornos concurrentes}: cada
  conexión al servidor mantiene su propio contador, por lo que dos sesiones
  simultáneas no interfieren entre sí.
\end{nota}

% ─────────────────────────────────────────────────────────
\subsection{Inserción en Tablas con Relación 1:N}

\begin{definicion}{Relación 1:N}
  Una fila de la tabla \textbf{padre} puede asociarse con \textbf{múltiples}
  filas de la tabla \textbf{hijo}, pero cada fila hijo pertenece a un único
  padre. Es la cardinalidad más frecuente en bases de datos relacionales.
\end{definicion}

\textbf{Esquema de ejemplo:}

\begin{lstlisting}[style=sql, caption={Tablas con relación 1:N — cliente y pedidos}]
CREATE TABLE clientes (
    cliente_id  INT          UNSIGNED AUTO_INCREMENT PRIMARY KEY,
    nombre      VARCHAR(100) NOT NULL,
    correo      VARCHAR(120) NOT NULL UNIQUE
);

CREATE TABLE pedidos (
    pedido_id   INT          UNSIGNED AUTO_INCREMENT PRIMARY KEY,
    cliente_id  INT          UNSIGNED NOT NULL,
    fecha       DATE         NOT NULL,
    total       DECIMAL(10,2) NOT NULL DEFAULT 0.00,
    FOREIGN KEY (cliente_id) REFERENCES clientes(cliente_id)
        ON DELETE RESTRICT
        ON UPDATE CASCADE
);
\end{lstlisting}

\textbf{Inserción de registros dependientes:}

\begin{lstlisting}[style=sql, caption={Inserción 1:N — un cliente con varios pedidos}]
-- Paso 1: insertar el cliente (tabla padre)
INSERT INTO clientes (nombre, correo)
VALUES ('Ana Torres', 'ana.torres@correo.com');

-- Paso 2: insertar sus pedidos (tabla hijo)
--         Todos referencian el mismo cliente_id
INSERT INTO pedidos (cliente_id, fecha, total)
VALUES
    (LAST_INSERT_ID(), '2024-01-15', 1250.00),
    (LAST_INSERT_ID(), '2024-03-02',  340.50),
    (LAST_INSERT_ID(), '2024-06-18', 2890.75);
\end{lstlisting}

\begin{ejemplo}
  Inserción de pedidos para un cliente ya existente, identificado por correo:

\begin{lstlisting}[style=sql, caption={INSERT ... SELECT para referenciar por atributo}]
INSERT INTO pedidos (cliente_id, fecha, total)
SELECT cliente_id, '2024-09-10', 780.00
FROM   clientes
WHERE  correo = 'ana.torres@correo.com';
\end{lstlisting}

  Este patrón evita tener que conocer el \texttt{cliente\_id} numérico de
  antemano y es útil en scripts de carga masiva.
\end{ejemplo}

\begin{nota}
  La política \texttt{ON DELETE RESTRICT} impide eliminar un cliente mientras
  tenga pedidos asociados. Para poder borrarlo habría que eliminar primero sus
  pedidos, o bien cambiar la política a \texttt{ON DELETE CASCADE}.
\end{nota}

% ─────────────────────────────────────────────────────────
\subsection{Inserción en Tablas con Relación N:M}

\begin{definicion}{Relación N:M}
  Múltiples filas de la tabla \textbf{A} pueden relacionarse con múltiples
  filas de la tabla \textbf{B}. Se implementa mediante una
  \textbf{tabla intermedia} (también llamada tabla de unión o pivote) que
  contiene, al menos, las claves foráneas de ambas entidades.
\end{definicion}

\textbf{Esquema de ejemplo:}

\begin{lstlisting}[style=sql, caption={Tablas con relación N:M — estudiantes y cursos}]
CREATE TABLE estudiantes (
    estudiante_id INT          UNSIGNED AUTO_INCREMENT PRIMARY KEY,
    nombre        VARCHAR(100) NOT NULL
);

CREATE TABLE cursos (
    curso_id      INT          UNSIGNED AUTO_INCREMENT PRIMARY KEY,
    titulo        VARCHAR(120) NOT NULL,
    creditos      TINYINT      UNSIGNED NOT NULL
);

-- Tabla intermedia: un estudiante puede tomar N cursos
--                  un curso puede tener N estudiantes
CREATE TABLE inscripciones (
    estudiante_id INT      UNSIGNED NOT NULL,
    curso_id      INT      UNSIGNED NOT NULL,
    fecha_alta    DATE     NOT NULL,
    calificacion  DECIMAL(4,2) DEFAULT NULL,
    PRIMARY KEY (estudiante_id, curso_id),    -- PK compuesta evita duplicados
    FOREIGN KEY (estudiante_id) REFERENCES estudiantes(estudiante_id)
        ON DELETE CASCADE ON UPDATE CASCADE,
    FOREIGN KEY (curso_id)      REFERENCES cursos(curso_id)
        ON DELETE CASCADE ON UPDATE CASCADE
);
\end{lstlisting}

\textbf{Orden de inserción en tres pasos:}

\begin{lstlisting}[style=sql, caption={Paso 1 — poblar ambas tablas padre}]
-- Tabla padre A
INSERT INTO estudiantes (nombre)
VALUES ('Miguel Ríos'),
       ('Sofía Castillo'),
       ('Héctor Núñez');

-- Tabla padre B
INSERT INTO cursos (titulo, creditos)
VALUES ('Bases de Datos',    6),
       ('Álgebra Lineal',    4),
       ('Redes de Cómputo',  5);
\end{lstlisting}

\begin{lstlisting}[style=sql, caption={Paso 2 — insertar en la tabla intermedia}]
-- IDs resultantes (ejemplo):
--   estudiante_id: 1=Miguel, 2=Sofía, 3=Héctor
--   curso_id:      1=Bases de Datos, 2=Álgebra, 3=Redes

INSERT INTO inscripciones (estudiante_id, curso_id, fecha_alta)
VALUES
    (1, 1, '2024-08-20'),   -- Miguel  -> Bases de Datos
    (1, 3, '2024-08-20'),   -- Miguel  -> Redes de Cómputo
    (2, 1, '2024-08-21'),   -- Sofía   -> Bases de Datos
    (2, 2, '2024-08-21'),   -- Sofía   -> Álgebra Lineal
    (3, 1, '2024-08-22'),   -- Héctor  -> Bases de Datos
    (3, 2, '2024-08-22'),   -- Héctor  -> Álgebra Lineal
    (3, 3, '2024-08-22');   -- Héctor  -> Redes de Cómputo
\end{lstlisting}

\begin{lstlisting}[style=sql, caption={Paso 3 — actualizar atributos propios de la relación}]
-- La calificacion es un atributo de la RELACIÓN, no de ninguna entidad sola
UPDATE inscripciones
SET    calificacion = 9.2
WHERE  estudiante_id = 1
  AND  curso_id      = 1;
\end{lstlisting}

\begin{ejemplo}
  Verificar las inscripciones de un estudiante con una consulta de control:

\begin{lstlisting}[style=sql, caption={Consulta de verificación post-inserción N:M}]
SELECT e.nombre        AS estudiante,
       c.titulo        AS curso,
       i.fecha_alta,
       i.calificacion
FROM   inscripciones AS i
INNER JOIN estudiantes AS e ON i.estudiante_id = e.estudiante_id
INNER JOIN cursos      AS c ON i.curso_id      = c.curso_id
WHERE  e.nombre = 'Miguel Ríos';
\end{lstlisting}
\end{ejemplo}

\begin{nota}
  La \textbf{clave primaria compuesta} \texttt{(estudiante\_id, curso\_id)} en
  la tabla intermedia garantiza que un mismo par no pueda insertarse dos veces,
  cumpliendo la regla de integridad de entidad sin necesidad de un campo
  \texttt{AUTO\_INCREMENT} adicional.
\end{nota}

% ── Resumen del módulo ────────────────────────────────────
\subsection*{Resumen del módulo}

\begin{table}[H]
  \centering
  \caption{Estrategias de inserción según la cardinalidad de la relación}
  \renewcommand{\arraystretch}{1.3}
  \rowcolors{2}{gray!10}{white}
  \begin{tabularx}{\textwidth}{l l X}
    \toprule
    \textbf{Relación} & \textbf{Tablas involucradas}     & \textbf{Consideración clave}                              \\
    \midrule
    1:1               & Padre + hijo con FK \texttt{UNIQUE} & Insertar padre primero; usar \texttt{LAST\_INSERT\_ID()}  \\
    1:N               & Padre + uno o más hijos          & Todos los hijos referencian el mismo padre                \\
    N:M               & Dos padres + tabla intermedia    & Insertar ambos padres antes de poblar la tabla pivote     \\
    \bottomrule
  \end{tabularx}
\end{table}